\subsubsection{Encoding Bit Strings into Quantum States}\label{sec:encoding-bit-strings-into-quantum-states}

\textcolor{Red2}{\faIcon{exclamation-triangle}} The idea of find the ground state of a Hamiltonian to solve a QUBO problem is powerful, but it raises an important question: \textbf{\emph{how do we encode the binary variables of the QUBO problem into quantum states?}}

\highspace
In a QUBO problem, every candidate solution is represented by a binary string:
\begin{equation*}
    \mathbf{x} = \left(x_{1}, x_{2}, \dots, x_{n}\right) \quad x_{i} \in \{0, 1\}
\end{equation*}
For example, the values $00$, $01$, $10$, and $11$ are possible solutions for a problem with two binary variables. However, a \hl{Hamiltonian acts on quantum states} rather than classical bit strings. Therefore, \hl{before defining the Hamiltonian, we must establish a correspondence between classical configurations and quantum states.}

\highspace
\begin{flushleft}
    \textcolor{Green3}{\faIcon{exchange-alt} \textbf{Mapping an Entire Bit String}}
\end{flushleft}
Let's start with the simplest case: a classical bit. It can take on two values, $0$ and $1$:
\begin{equation*}
    x_{i} \in \{0, 1\}
\end{equation*}
It is simply mapped to a quantum state of a single qubit:
\begin{equation*}
    \ket{0} = \begin{bmatrix}
        1 \\
        0
    \end{bmatrix} \leftrightarrow 0, \quad \ket{1} = \begin{bmatrix}
        0 \\
        1
    \end{bmatrix} \leftrightarrow 1
\end{equation*}
This is simply the standard computational basis of quantum computing.

\highspace
In general, a QUBO solution contains $n$ bits:
\begin{equation*}
    \mathbf{x} = \left(x_1, x_2, \dots, x_n\right)
\end{equation*}
Each bit is mapped to a qubit. The entire string is then represented using the tensor product of the individual qubits:
\begin{equation*}
    \ket{\mathbf{x}} = \ket{x_1} \otimes \ket{x_2} \otimes \cdots \otimes \ket{x_n}
\end{equation*}
Therefore, in general a bit string of length $n$ corresponds to a quantum state in a $2^n$-dimensional Hilbert space.

\begin{examplebox}[: Encoding the Bit String $10$]
    Consider the bit string $x = \texttt{10}$.
    \begin{equation*}
        1 \iff \ket{1} = \begin{bmatrix}
            0 \\
            1
        \end{bmatrix}
        \quad
        0 \iff \ket{0} = \begin{bmatrix}
            1 \\
            0
        \end{bmatrix}
    \end{equation*}
    The corresponding two-qubit state is:
    \begin{equation*}
        \ket{10} = \ket{1} \otimes \ket{0} = \begin{bmatrix}
            0 \\
            1
        \end{bmatrix} \otimes \begin{bmatrix}
            1 \\
            0
        \end{bmatrix} = \begin{bmatrix}
            0 \\
            0 \\
            1 \\
            0
        \end{bmatrix}
    \end{equation*}
\end{examplebox}